\documentclass[11pt,a4paper]{article}

%, Packages,
\usepackage[utf8]{inputenc}
\usepackage[T1]{fontenc}
\usepackage[french]{babel}
\usepackage{lmodern}
\usepackage[margin=2.5cm]{geometry}
\usepackage{amsmath,amssymb}
\usepackage{booktabs}
\usepackage{array}
\usepackage{enumitem}
\usepackage{hyperref}
\usepackage{xcolor}
\usepackage{parskip}

\hypersetup{
    colorlinks=true,
    linkcolor=blue!60!black,
    urlcolor=blue!60!black,
    citecolor=blue!60!black
}

%, Métadonnées,
\title{Résumé simple du papier \\
\large \textit{Slow Momentum with Fast Reversion: A Trading Strategy Using Deep Learning and Changepoint Detection}}
\author{Wood, K., Roberts, S., Zohren, S. (2022) \\ \small The Journal of Financial Data Science}
\date{Projet Pergam MSc 2026 -- ESILV}

\begin{document}

\maketitle

\section{Le problème de départ}

Les stratégies de \textbf{momentum} reposent sur une idée simple : \emph{« une tendance qui monte continue de monter, une tendance qui baisse continue de baisser »}. On suit donc la tendance (on achète ce qui monte, on vend ce qui baisse).

Ça marche bien... \textbf{sauf aux moments où la tendance se retourne brutalement} (ce que les auteurs appellent \emph{momentum turning points}). Exemple typique : le crash COVID de mars 2020. Le marché montait, puis s'effondre en quelques jours. Une stratégie momentum classique est encore « long » (acheteuse) au mauvais moment et prend la chute de plein fouet.

Les auteurs constatent que \textbf{même les stratégies récentes basées sur du deep learning} (les \emph{Deep Momentum Networks} de Lim, Zohren \& Roberts 2019) souffrent du même problème sur la période 2015--2020.

\section{L'idée centrale du papier}

Combiner \textbf{trois briques} :
\begin{enumerate}[leftmargin=*]
    \item \textbf{Slow momentum} (momentum lent) $\rightarrow$ suit les vraies tendances de fond (ex : sur 1 an).
    \item \textbf{Fast mean-reversion} (retour à la moyenne rapide) $\rightarrow$ parie que les petits mouvements court terme vont se corriger (ex : sur 1 mois).
    \item \textbf{Un détecteur de changements de régime} (\emph{Changepoint Detection}, CPD) $\rightarrow$ dit au modèle « attention, quelque chose vient de changer dans le marché ».
\end{enumerate}

Le tout est branché dans un \textbf{réseau de neurones LSTM} qui apprend tout seul à doser les deux signaux en fonction du contexte.

\begin{quote}
\textit{Image mentale :} c'est comme un conducteur qui roule vite sur autoroute (slow momentum), mais qui a un copilote (le CPD) qui lui dit « freine, il y a un virage » dès qu'un changement de régime arrive.
\end{quote}

\section{Les grandes parties du papier}

\subsection{Détection de changements de régime (Changepoint Detection)}

\textbf{But :} détecter automatiquement quand le comportement du marché change.

\textbf{Comment :}
\begin{itemize}[leftmargin=*]
    \item On prend les rendements journaliers sur une fenêtre de $l$ jours (appelée LBW = \emph{lookback window}).
    \item On ajuste dessus un \textbf{Processus Gaussien} (GP), un modèle statistique bayésien qui « dessine » une courbe probable à travers les données.
    \item On utilise un noyau \textbf{Matérn 3/2} (adapté aux données bruitées et non lisses comme la finance).
    \item Ensuite, on compare deux hypothèses :
    \begin{itemize}
        \item \textbf{H1} : les données suivent un seul régime (un seul noyau).
        \item \textbf{H2} : il y a un changement à un instant $c$, avec deux noyaux différents (un avant, un après), reliés par une transition en sigmoïde.
    \end{itemize}
    \item Si H2 explique \textbf{beaucoup mieux} les données que H1, c'est qu'il y a eu un vrai changement.
\end{itemize}

\textbf{Ce qu'on en sort :}
\begin{itemize}[leftmargin=*]
    \item $\nu$ (la \textbf{sévérité}) : entre 0 et 1, indique la force du changement.
    \item $\gamma$ (la \textbf{localisation}) : entre 0 et 1, indique \emph{où} dans la fenêtre s'est produit le changement (proche $=$ récent, loin $=$ vieux).
\end{itemize}

Ces deux chiffres deviennent des \textbf{entrées supplémentaires} pour le réseau de neurones.

\subsection{Deep Momentum Network (LSTM)}

\textbf{But :} apprendre directement la position à prendre (entre $-1$ $=$ vente max et $+1$ $=$ achat max).

\textbf{Comment :}
\begin{itemize}[leftmargin=*]
    \item Architecture $=$ \textbf{LSTM} (Long Short-Term Memory), un type de réseau récurrent bon pour les séries temporelles.
    \item \textbf{Fonction de perte $=$ Sharpe ratio négatif}. Autrement dit, le réseau n'essaie pas de prédire si le prix monte ou descend, il essaie directement de \textbf{maximiser le rendement ajusté du risque}. C'est une astuce très importante : prédire la direction n'est pas la même chose que gagner de l'argent.
    \item \textbf{Entrées du modèle} :
    \begin{itemize}
        \item Rendements normalisés sur plusieurs horizons (1, 21, 63, 126, 252 jours).
        \item Indicateurs MACD (3 paires de moyennes mobiles).
        \item \textbf{Sévérité $\nu$ et localisation $\gamma$} issues du CPD.
    \end{itemize}
    \item \textbf{Sortie} : une position $X$ entre $-1$ et $1$, obtenue via une fonction $\tanh$.
\end{itemize}

\textbf{Scaling en volatilité :} avant tout, on ajuste la taille des positions pour que chaque actif contribue autant au risque total (cible $=$ 15\% de volatilité annualisée).

\subsection{Stratégie et backtest}

\textbf{Univers :} 50 contrats futures liquides (matières premières, actions, taux, devises), de 1990 à 2020, données Pinnacle.

\textbf{Méthode de test (expanding window) :}
\begin{itemize}[leftmargin=*]
    \item Entraîne sur 1990--1995 $\rightarrow$ teste sur 1995--2000.
    \item Étend à 1990--2000 $\rightarrow$ teste sur 2000--2005.
    \item Et ainsi de suite jusqu'en 2020.
    \item À chaque étape, les hyperparamètres sont ré-optimisés.
\end{itemize}

\textbf{Benchmarks comparés :}
\begin{itemize}[leftmargin=*]
    \item Long Only (acheter et garder)
    \item MACD classique
    \item TSMOM classique (Moskowitz 2012) avec différents mélanges slow/fast
    \item LSTM seul (sans CPD)
    \item \textbf{LSTM $+$ CPD} (la contribution du papier)
\end{itemize}

\section{Résultats principaux}

Sur la période complète (1995--2020), rescalé à 15\% de volatilité annualisée :

\begin{center}
\begin{tabular}{lrr}
\toprule
\textbf{Stratégie} & \textbf{Sharpe} & \textbf{Rendement} \\
\midrule
Long Only            & 0.44 & 6.6\% \\
MACD                 & 0.77 & 11.1\% \\
TSMOM classique      & 0.94 & 13.8\% \\
LSTM seul            & 1.62 & 21.0\% \\
\textbf{LSTM $+$ CPD} & \textbf{2.16} & \textbf{31.5\%} \\
\bottomrule
\end{tabular}
\end{center}

\textbf{Gains clés :}
\begin{itemize}[leftmargin=*]
    \item \textbf{$+33\%$ de Sharpe} vs. LSTM sans CPD sur toute la période.
    \item \textbf{$+66\%$ de Sharpe} sur 2015--2020 (la période la plus turbulente, où les stratégies momentum classiques souffrent le plus).
    \item Même une petite fenêtre de 2 semaines pour le CPD apporte déjà un gain notable.
    \item La meilleure fenêtre fixe est $\sim$1 trimestre (63 jours), mais laisser le LBW comme hyperparamètre optimisé donne le meilleur résultat.
\end{itemize}

\section{Les limites du papier}

\subsection{Pas de coûts de transaction}
C'est \textbf{la} limite la plus importante. Le papier reconnaît qu'il ne modélise pas les frais. Or la stratégie fait du \emph{fast mean-reversion}, donc \textbf{elle tourne beaucoup} (elle retourne souvent sa position). En conditions réelles, les coûts peuvent manger une grosse partie du Sharpe. Un appendice donne quelques éléments, mais ce n'est pas intégré à l'entraînement.

\subsection{Coût de calcul}
Ajuster un GP avec noyau de changepoint \textbf{pour chaque actif, chaque jour, et chaque taille de fenêtre} est lourd. Pour de grandes fenêtres (126, 252 jours), ça devient vite prohibitif. En pratique, ça limite la fréquence à laquelle on peut ré-entraîner.

\subsection{Choix arbitraires dans le CPD}
\begin{itemize}[leftmargin=*]
    \item Le noyau Matérn 3/2 est choisi empiriquement.
    \item La fenêtre LBW est sensible : trop courte $=$ bruit, trop longue $=$ on rate les changements.
    \item L'hypothèse d'\textbf{un seul changepoint} dans la fenêtre est forte, un vrai marché peut avoir plusieurs ruptures rapprochées.
\end{itemize}

\subsection{Robustesse et overfitting}
\begin{itemize}[leftmargin=*]
    \item Beaucoup d'hyperparamètres (LSTM $+$ GP $+$ LBW), optimisés sur validation : \textbf{risque de sur-ajustement} aux fenêtres historiques.
    \item Le papier ne teste qu'\textbf{un seul tirage} de la procédure d'entraînement : pas d'intervalle de confiance, pas de test statistique formel sur les différences de Sharpe.
\end{itemize}

\subsection{Univers limité}
50 futures seulement, tous liquides. Pas d'actions individuelles, pas de crypto, pas de marchés émergents. La généralisation à d'autres classes d'actifs n'est pas prouvée.

\subsection{Pas de contrainte de portefeuille réaliste}
Pas de limite de levier, pas de contrainte de concentration, pas de gestion du risque globale (VaR, drawdown cap). Juste 50 positions indépendantes additionnées.

\subsection{Stationnarité résiduelle}
Même avec le CPD, le modèle suppose que les relations apprises par le LSTM restent valides. Sur un changement structurel majeur (nouveau régime monétaire, guerre, pandémie inédite), rien ne garantit que ça tienne.

\subsection{Boîte noire}
Le LSTM reste \textbf{peu interprétable}. Difficile d'expliquer \emph{pourquoi} une position est prise à un instant donné, problème classique pour la conformité et la confiance d'un gestionnaire de fonds.

\section{Pistes d'amélioration intéressantes (pour le projet)}

\begin{enumerate}[leftmargin=*]
    \item \textbf{Intégrer les coûts de transaction directement dans la loss Sharpe} $\rightarrow$ c'est le plus important et le plus réaliste.
    \item \textbf{Tester d'autres méthodes de CPD} : Bayesian Online Changepoint Detection (BOCPD), \texttt{ruptures}, CUSUM, ou même un petit réseau de neurones dédié.
    \item \textbf{Remplacer le LSTM par un Transformer} (Temporal Fusion Transformer) pour voir si l'attention aide.
    \item \textbf{Utiliser plusieurs LBW en parallèle} et laisser le réseau choisir (les auteurs ont essayé et ça n'a pas marché, mais avec une autre architecture ça pourrait).
    \item \textbf{Étendre au post-2020} pour voir si la stratégie tient toujours (inflation 2022, taux qui remontent\ldots).
    \item \textbf{Étudier la stabilité} : relancer l'entraînement avec plusieurs seeds et regarder la variance des résultats.
    \item \textbf{Ajouter un risk management global} : cap de drawdown, de levier, corrélations entre actifs.
\end{enumerate}

\end{document}
